\documentclass[10pt,a4paper]{article}
\usepackage{nexusS5}
\setnexusheader{Dossier enseignant — confidentiel}{S5 — Malek KHADHRANI — usage enseignant}
\begin{document}
\nexuscover{SÉANCE 5 --- DOSSIER ENSEIGNANT}{Profil, déroulé, corrigés et barème}{Malek KHADHRANI}{Entrée en Première générale — Spécialité mathématiques}{Mathématiques --- Substitution rigoureuse et vecteurs}
\begin{methodbox}\small\textbf{Programme de référence.} nouveau programme de la spécialité mathématiques de Première --- BO n° 14 du 2 avril 2026 MENE2602917A. Transition visée : acquis de Seconde 2025-2026 vers le nouveau programme de spécialité de Première applicable en 2026-2027.\end{methodbox}
\begin{keybox}\textbf{DOCUMENT CONFIDENTIEL --- USAGE ENSEIGNANT.} Contient les corrigés, le barème et des données nominatives. Ne pas remettre à l'élève ni le laisser visible pendant la séance.\end{keybox}
\sectiontitle{1. Profil synthétique}
\begin{methodbox}
\textbf{Diagnostic initial.} Diagnostic initial (18 items). Calcul littéral et pourcentages solides ; fonctions, fonctions de référence, vecteurs et calcul numérique à rectifier ; droites à installer.
\end{methodbox}
\subsectiontitle{Points d'appui documentés}
\begin{itemize}\item Calcul littéral\item Pourcentages et évolutions\end{itemize}
\subsectiontitle{Difficultés documentées}
\begin{itemize}\item Substitution négative\item Image/antécédent\item Fonction inverse\item Coordonnées de vecteur\item Colinéarité\item Pente\item Puissances\end{itemize}
\begin{warningbox}\textbf{Limite de la reconstruction.} Les tableaux d'observation séance par séance du dossier individuel sont vierges : il n'existe aucune preuve documentée entre le diagnostic initial et aujourd'hui. Tout statut porté ci-dessous décrit un \emph{état de départ}, jamais une acquisition constatée.\end{warningbox}
\newpage\sectiontitle{2. Matrice de trajectoire --- état avant la séance 5}
\rowcolors{2}{SoftBlue}{white}
\par\noindent\begin{tabularx}{\linewidth}{>{\ttfamily\scriptsize}p{27mm} X >{\scriptsize}p{16mm} >{\scriptsize}p{20mm} >{\scriptsize\centering\arraybackslash}p{9mm}}
\rowcolor{Navy}\color{white}\scriptsize skill\_id & \color{white}Compétence & \color{white}Travaillée en & \color{white}Statut avant S5 & \color{white}Prio. \\
M1RE\_ALG\_01 & Développement et identités remarquables & S1 & acquis & OK \\
M1RE\_ALG\_02 & Factorisation, en particulier la différence de deux carrés & S1 & acquis & OK \\
M1RE\_POURC\_01 & Taux d'évolution & S4 & acquis & OK \\
M1RE\_POURC\_02 & Coefficient multiplicateur et évolutions successives & S4 & acquis & OK \\
M1RE\_FONC\_01 & $\bullet$ Image d'un réel, substitution d'une valeur négative & S2 & fragile & P1 \\
M1RE\_FONC\_02 & Langage des fonctions : image, antécédent, appartenance à la courbe & S2 & fragile & P1 \\
M1RE\_VECT\_01 & Coordonnées d'un vecteur défini par deux points & S3 & fragile & P1 \\
M1RE\_VECT\_02 & $\bullet$ Colinéarité de deux vecteurs & S3 & fragile & P1 \\
M1RE\_DROIT\_01 & Coefficient directeur d'une droite & S3 & en voie & P2 \\
M1RE\_FREF\_01 & Fonctions de référence : carré et inverse, variations & S2 & fragile & P2 \\
M1RE\_FREF\_02 & Comparaison de x$^2$ et de x selon la position de x & S2 & fragile & P2 \\
M1RE\_INEQ\_01 & Inéquation du premier degré et écriture en intervalle & S1 & en voie & P2 \\
M1RE\_INEQ\_02 & Signe d'un produit et tableau de signes & S1 & en voie & P2 \\
M1RE\_NUM\_01 & Calcul avec des puissances & S1 & fragile & P2 \\
M1RE\_NUM\_02 & Simplification d'un radical & S1 & fragile & P2 \\
M1RE\_DROIT\_02 & Lecture d'une équation réduite de droite & S3 & en voie & P3 \\
M1RE\_PROBA\_01 & Union, intersection et formule P(A$\cup$B) & S4 & en voie & P3 \\
M1RE\_PROBA\_02 & Événement contraire & S4 & en voie & P3 \\
\end{tabularx}
\par\vspace{1mm}\small\textbf{Lecture de la colonne « Statut avant S5 » :} elle décrit l'état constaté au \emph{positionnement initial}, jamais une acquisition constatée depuis. « acquis » signifie « acquis au positionnement initial » ; « en voie » signifie « en voie d'acquisition ».
\par\small$\bullet$ compétence ciblée par les phases 2 et 3 de la séance. La colonne « Prio. » est \emph{provisoire} : la priorité définitive est produite par \code{tools/analyze_s5.py} après la correction.
\begin{warningbox}\textbf{Effet de récence.} Les compétences suivantes sont retravaillées pendant cette séance, puis évaluées moins d'une heure plus tard : \code{M1RE_FONC_01}, \code{M1RE_VECT_02}. Une réussite y mesure une \emph{performance immédiate}, pas une consolidation durable. Le plan de rentrée prévoit pour elles un contrôle différé en semaine 1 ou 2 ; aucun bilan ne doit écrire « acquis durablement » sur cette seule preuve.\end{warningbox}
\sectiontitle{3. Pourquoi ces activités, pour cet élève}
\begin{goalbox}\textbf{Objectif de la séance :} Ralentir lorsque la question paraît familière et contrôler la règle activée — conseil du dossier — en sécurisant la substitution d'une valeur négative puis le test de colinéarité.\end{goalbox}
\subsectiontitle{Axe prioritaire (phase 2, 25 min) --- Calculer une image, en particulier pour une valeur négative}
Erreur documentée sur la substitution d'un nombre négatif ; l'omission des parenthèses est une erreur d'écriture qui produit systématiquement un résultat faux.
\par\vspace{1mm}\textbf{Compétence visée :} \code{M1RE_FONC_01}.
\subsectiontitle{Second axe (phase 3, 20 min) --- Colinéarité de deux vecteurs et alignement de trois points}
La colinéarité figure parmi les priorités du dossier ; c'est le critère utilisé toute l'année pour l'alignement, le parallélisme et l'équation de droite.
\par\vspace{1mm}\textbf{Compétence visée :} \code{M1RE_VECT_02}.
\subsectiontitle{Équilibre commun / individuel}
Sur les 75 minutes : \textbf{51 min} (68\,\%) relèvent des objectifs communs du niveau et \textbf{24 min} (32\,\%) ciblent le profil individuel. Sur l'évaluation : \textbf{14 points} de noyau commun (70.0\,\%) et \textbf{6 points} individualisés (30.0\,\%).
\newpage\sectiontitle{4. Déroulé minute par minute}
\rowcolors{2}{SoftGray}{white}
\par\noindent\begin{tabularx}{\linewidth}{>{\bfseries}p{18mm} p{34mm} X X}
\rowcolor{Navy}\color{white}Temps & \color{white}Phase & \color{white}Conduite & \color{white}Indicateur observable \\
0--10 min & Réactiver les cinq blocs du stage (algèbre, fonctions, vecte & Individuel, sans carte de rappel pendant 7 minutes, puis 3 minutes de mise en commun. & Au moins 4 réponses exactes sur 5 ; la certitude est renseignée avant toute vérification. \\
10--35 min & Calculer une image, en particulier pour une valeur négative & Rappel bref, exemple guidé au tableau, puis exercices en autonomie. Aide donnée seulement après une tentative écrite. & Trois réussites consécutives sans aide sur la compétence visée. \\
35--55 min & Colinéarité de deux vecteurs et alignement de trois points & Pas d'exemple guidé : l'élève démarre seul. Intervention uniquement sur demande. & Deux réussites autonomes ; le contrôle est écrit sans qu'on le demande. \\
55--70 min & De la Seconde à la Première : la suite et le taux de variati & Travail écrit, correction collective courte à la fin. Ne pas transformer ce temps en cours magistral. & L'élève relie explicitement un acquis de l'année écoulée à un attendu de l'année à venir. \\
70--75 min & Cadrage méthodologique, sans aucune information sur le conte & Lecture des cinq règles, puis remise de l'évaluation. Aucune information sur son contenu. & Les deux engagements sont écrits. \\
75--120 min & Évaluation finale & Individuel, sans document. Partie A environ 10 min, B environ 20 min, C environ 15 min. & Somme des durées cibles : 37 min, marge de 7 min. \\
\end{tabularx}
\newpage\sectiontitle{5. Corrigé du livret de travail}
\subsectiontitle{Phase 1 --- réactivation}
\begin{enumerate}
\item \code{M1RE_ALG_02} --- $(3x)^2-5^2=\mathbf{(3x-5)(3x+5)}$. Contrôle : développer redonne $9x^2-25$.
\item \code{M1RE_INEQ_01} --- $-4x\ge -10$ donc $x\le \frac52$ : $S=\left]-\infty\,;\tfrac52\right]$ (sens de l'inégalité changé).
\item \code{M1RE_FONC_01} --- $f(-3)=(-3)^2-5\times(-3)=9+15=\mathbf{24}$.
\item \code{M1RE_VECT_01} --- $\vec{AB}(3-(-2)\,;-1-5)=\mathbf{(5\,;-6)}$ ; coefficient directeur $=\frac{-6}{5}=\mathbf{-1{,}2}$.
\item \code{M1RE_POURC_01} --- $\frac{210-250}{250}=-0{,}16$, soit une baisse de $\mathbf{16\,\%}$.
\end{enumerate}
\subsectiontitle{Phase 2 --- Calculer une image, en particulier pour une valeur négative}
\begin{enumerate}
\item $f(-4)=16-8=\mathbf{8}$ ; $f(0)=\mathbf{0}$ ; $f(3)=9+6=\mathbf{15}$.
\item $g(-3)=-2\times 9 + (-3) = \mathbf{-21}$.
\item $h(-4)=(-5)^2=\mathbf{25}$.
\item $25$ et $-25$ : dans le second cas, le signe $-$ est appliqué après le carré.
\item $k(-2)=-8$ : un exposant impair conserve le signe du nombre.
\item Les parenthèses manquent autour de $-2$ : il fallait $(-2)^2 = 4$. Résultat exact : $f(-2)=4+6=\mathbf{10}$. Contrôle : les deux termes sont positifs, le résultat ne peut valoir $2$.
\end{enumerate}
\minititle{Erreurs probables}
\begin{itemize}
\item \textsc{NOTATION} --- parenthèses omises autour d'une valeur négative substituée
\item \textsc{CALCUL} --- signe du produit $-4\times(-3)$ erroné
\item \textsc{CONTROLE} --- signe du résultat non confronté aux signes des termes
\end{itemize}
\subsectiontitle{Phase 3 --- Colinéarité de deux vecteurs et alignement de trois points}
\begin{enumerate}
\item $5\times 6 - (-2)\times(-15) = 30 - 30 = 0$ : ils sont \textbf{colinéaires}.
\item $4\times 3 - 7\times 2 = 12 - 14 = -2 \neq 0$ : ils ne sont \textbf{pas} colinéaires.
\item $\vec{AB}(3\,;6)$ et $\vec{AC}(-2\,;-4)$ ; déterminant $3\times(-4)-6\times(-2)=-12+12=0$ : les points sont alignés.
\item $3\times 15 - k\times 9 = 0$ donne $45 = 9k$ donc $k = \mathbf{5}$.
\end{enumerate}
\minititle{Erreurs probables}
\begin{itemize}
\item \textsc{METHODE} --- conclusion tirée d'une seule coordonnée
\item \textsc{CALCUL} --- signe perdu dans le produit $-y\times x'$
\item \textsc{JUSTIFICATION} --- colinéarité affirmée sans calcul du déterminant
\end{itemize}
\subsectiontitle{Phase 4 --- pont vers Première spécialité}

\textbf{A.} 1. Coefficient $1{,}03$ ; $u_0=2000$. \quad 2. $u_{n+1}=1{,}03\,u_n$.
3. $u_1=2060$ ; $u_2=2121{,}80$. \quad 4. \code{u = 1.03 * u} puis \code{print(n + 1, u)}.
La ligne \code{print(0, u)} placée \emph{avant} la boucle affiche $u_0$ ; sans elle, le programme
afficherait seulement $u_1$ à $u_5$, car l'affectation précède l'affichage à chaque tour.
5. $u_3=2000\times 1{,}03^3\approx 2185{,}45 < 2200$ : non, le seuil n'est pas atteint en trois ans.

\textbf{B.} 1. $\frac{f(3)-f(1)}{3-1}=\frac{9-1}{2}=4$.
2. $\frac{2{,}25-1}{0{,}5}=2{,}5$ ; $\frac{1{,}21-1}{0{,}1}=2{,}1$.
3. Les taux se rapprochent de $2$.
4. Attendu : ce nombre décrira la pente de la tangente à la courbe au point d'abscisse $1$ ; c'est le nombre dérivé.

\textbf{Observation à noter :} écrire $u_{n+1}=u_n+3$ est une erreur de \textsc{CONCEPT} (évolution additive
au lieu d'une évolution multiplicative) ; inverser numérateur et dénominateur du taux de variation est une
erreur de \textsc{METHODE}.

\newpage\sectiontitle{6. Corrigé de l'évaluation et barème analytique}
\begin{keybox}\textbf{Total : 20 points sur 20.} Noyau commun 14 pts (70.0\,\%), items individualisés 6 pts (30.0\,\%). Somme des durées cibles : 37 minutes.\end{keybox}
\itemhead{A1 --- \code{M1RE_ALG_02} --- noyau commun}{1}{1}
\begin{graybox}\small\textbf{Réponse attendue :} $(4x - 7)(4x + 7)$\end{graybox}
\minititle{Étapes significatives}
\begin{enumerate}\item reconnaître $a^2-b^2$ avec $a=4x$ et $b=7$\end{enumerate}
\minititle{Barème par critère --- la saisie se fait critère par critère}
\par\noindent\begin{tabularx}{\linewidth}{>{\ttfamily\scriptsize}p{9mm} >{\bfseries}p{9mm} >{\ttfamily\scriptsize}p{24mm} >{\scriptsize}p{18mm} X}
\toprule \scriptsize critère & Pts & \scriptsize compétence & \scriptsize preuve & Ce qui est exigé \\ \midrule
c1 & 1 & M1RE\_ALG\_02 & procedure & factorisation exacte \\
\bottomrule\end{tabularx}
\minititle{Erreurs probables et codes}
\par\noindent\begin{tabularx}{\linewidth}{>{\ttfamily\small}p{26mm} X}
\toprule Code & Description \\ \midrule
CONCEPT & $(4x-7)^2$ : différence de carrés confondue avec un carré parfait \\
CALCUL & racine de $16x^2$ prise égale à $16x$ \\
\bottomrule\end{tabularx}
\par\vspace{1mm}\small\textbf{Rôle pour la rentrée :} outil direct du second degré et des études de signe
\par\vspace{2mm}
\itemhead{A2 --- \code{M1RE_INEQ_01} --- noyau commun}{1}{1}
\begin{graybox}\small\textbf{Réponse attendue :} $x < 4$, soit $S = \left]-\infty\,;4\right[$\end{graybox}
\minititle{Étapes significatives}
\begin{enumerate}\item $-3x > -12$\item division par $-3$ : le sens de l'inégalité change\item $x<4$\end{enumerate}
\minititle{Barème par critère --- la saisie se fait critère par critère}
\par\noindent\begin{tabularx}{\linewidth}{>{\ttfamily\scriptsize}p{9mm} >{\bfseries}p{9mm} >{\ttfamily\scriptsize}p{24mm} >{\scriptsize}p{18mm} X}
\toprule \scriptsize critère & Pts & \scriptsize compétence & \scriptsize preuve & Ce qui est exigé \\ \midrule
c1 & 1 & M1RE\_INEQ\_01 & procedure & solution et intervalle exacts \\
\bottomrule\end{tabularx}
\minititle{Erreurs probables et codes}
\par\noindent\begin{tabularx}{\linewidth}{>{\ttfamily\small}p{26mm} X}
\toprule Code & Description \\ \midrule
CONCEPT & sens de l'inégalité conservé après division par un négatif \\
NOTATION & crochet fermé sur une inégalité stricte \\
\bottomrule\end{tabularx}
\par\vspace{1mm}\small\textbf{Rôle pour la rentrée :} prérequis des tableaux de signes du second degré
\par\vspace{2mm}
\itemhead{A3 --- \code{M1RE_FONC_01} --- noyau commun}{1}{1}
\begin{graybox}\small\textbf{Réponse attendue :} $15$\end{graybox}
\minititle{Étapes significatives}
\begin{enumerate}\item $g(-3)=2\times(-3)^2+(-3)$\item $=2\times 9-3=15$\end{enumerate}
\minititle{Barème par critère --- la saisie se fait critère par critère}
\par\noindent\begin{tabularx}{\linewidth}{>{\ttfamily\scriptsize}p{9mm} >{\bfseries}p{9mm} >{\ttfamily\scriptsize}p{24mm} >{\scriptsize}p{18mm} X}
\toprule \scriptsize critère & Pts & \scriptsize compétence & \scriptsize preuve & Ce qui est exigé \\ \midrule
c1 & 1 & M1RE\_FONC\_01 & procedure & $15$ exact avec la substitution visible \\
\bottomrule\end{tabularx}
\minititle{Erreurs probables et codes}
\par\noindent\begin{tabularx}{\linewidth}{>{\ttfamily\small}p{26mm} X}
\toprule Code & Description \\ \midrule
NOTATION & parenthèses omises : $-3^2$ traité comme $-9$ \\
CALCUL & $2\times 9$ mal évalué \\
\bottomrule\end{tabularx}
\par\vspace{1mm}\small\textbf{Rôle pour la rentrée :} substitution employée en permanence en dérivation et sur les suites
\par\vspace{2mm}
\itemhead{A4 --- \code{M1RE_VECT_01} --- noyau commun}{1}{1}
\begin{graybox}\small\textbf{Réponse attendue :} $\vec{AB}(-6\,;6)$\end{graybox}
\minititle{Étapes significatives}
\begin{enumerate}\item $x_B-x_A=-2-4=-6$\item $y_B-y_A=5-(-1)=6$\end{enumerate}
\minititle{Barème par critère --- la saisie se fait critère par critère}
\par\noindent\begin{tabularx}{\linewidth}{>{\ttfamily\scriptsize}p{9mm} >{\bfseries}p{9mm} >{\ttfamily\scriptsize}p{24mm} >{\scriptsize}p{18mm} X}
\toprule \scriptsize critère & Pts & \scriptsize compétence & \scriptsize preuve & Ce qui est exigé \\ \midrule
c1 & 1 & M1RE\_VECT\_01 & automatisme & les deux coordonnées exactes \\
\bottomrule\end{tabularx}
\minititle{Erreurs probables et codes}
\par\noindent\begin{tabularx}{\linewidth}{>{\ttfamily\small}p{26mm} X}
\toprule Code & Description \\ \midrule
CONCEPT & soustraction effectuée dans l'ordre $A-B$ \\
CALCUL & $5-(-1)$ traité comme $4$ \\
\bottomrule\end{tabularx}
\par\vspace{1mm}\small\textbf{Rôle pour la rentrée :} prérequis du produit scalaire et de la géométrie repérée
\par\vspace{2mm}
\itemhead{A5 --- \code{M1RE_FONC_01} --- item individualisé}{1}{1.25}
\begin{graybox}\small\textbf{Réponse attendue :} $40$\end{graybox}
\minititle{Étapes significatives}
\begin{enumerate}\item $f(-4) = (-4)^2 - 6\times(-4)$\item $= 16 + 24 = 40$\end{enumerate}
\minititle{Barème par critère --- la saisie se fait critère par critère}
\par\noindent\begin{tabularx}{\linewidth}{>{\ttfamily\scriptsize}p{9mm} >{\bfseries}p{9mm} >{\ttfamily\scriptsize}p{24mm} >{\scriptsize}p{18mm} X}
\toprule \scriptsize critère & Pts & \scriptsize compétence & \scriptsize preuve & Ce qui est exigé \\ \midrule
c1 & 1 & M1RE\_FONC\_01 & procedure & $40$ exact avec la substitution visible \\
\bottomrule\end{tabularx}
\minititle{Erreurs probables et codes}
\par\noindent\begin{tabularx}{\linewidth}{>{\ttfamily\small}p{26mm} X}
\toprule Code & Description \\ \midrule
NOTATION & parenthèses omises : $-4^2$ traité comme $-16$ \\
CALCUL & signe du produit $-6\times(-4)$ erroné \\
\bottomrule\end{tabularx}
\par\vspace{1mm}\small\textbf{Rôle pour la rentrée :} substitution employée en dérivation et sur les suites
\par\vspace{2mm}
\itemhead{A6 --- \code{M1RE_VECT_02} --- item individualisé}{1}{1.5}
\begin{graybox}\small\textbf{Réponse attendue :} $6\times 6 - (-4)\times(-9) = 36 - 36 = 0$ : ils sont colinéaires.\end{graybox}
\minititle{Étapes significatives}
\begin{enumerate}\item déterminant $xy' - yx'$\end{enumerate}
\minititle{Barème par critère --- la saisie se fait critère par critère}
\par\noindent\begin{tabularx}{\linewidth}{>{\ttfamily\scriptsize}p{9mm} >{\bfseries}p{9mm} >{\ttfamily\scriptsize}p{24mm} >{\scriptsize}p{18mm} X}
\toprule \scriptsize critère & Pts & \scriptsize compétence & \scriptsize preuve & Ce qui est exigé \\ \midrule
c1 & 1 & M1RE\_VECT\_02 & communication & déterminant calculé et conclusion exacte \\
\bottomrule\end{tabularx}
\minititle{Erreurs probables et codes}
\par\noindent\begin{tabularx}{\linewidth}{>{\ttfamily\small}p{26mm} X}
\toprule Code & Description \\ \midrule
CALCUL & signe perdu dans $-y\times x'$ \\
JUSTIFICATION & conclusion sans calcul \\
\bottomrule\end{tabularx}
\par\vspace{1mm}\small\textbf{Rôle pour la rentrée :} critère d'alignement et de parallélisme, utilisé toute l'année
\par\vspace{2mm}
\itemhead{B1 --- \code{M1RE_INEQ_02} --- noyau commun}{2}{5.25}
\begin{graybox}\small\textbf{Réponse attendue :} Racines $-2$ et $2$ ; $P(x)>0$ sur $\left]-2\,;2\right[$ ; $P(x)\leqslant 0$ sur $\left]-\infty\,;-2\right]\cup\left[2\,;+\infty\right[$.\end{graybox}
\minititle{Étapes significatives}
\begin{enumerate}\item $3x+6=0 \Leftrightarrow x=-2$ ; $2-x=0 \Leftrightarrow x=2$\item $3x+6$ négatif avant $-2$, positif après ; $2-x$ positif avant $2$, négatif après\item produit positif entre les deux racines\end{enumerate}
\minititle{Barème par critère --- la saisie se fait critère par critère}
\par\noindent\begin{tabularx}{\linewidth}{>{\ttfamily\scriptsize}p{9mm} >{\bfseries}p{9mm} >{\ttfamily\scriptsize}p{24mm} >{\scriptsize}p{18mm} X}
\toprule \scriptsize critère & Pts & \scriptsize compétence & \scriptsize preuve & Ce qui est exigé \\ \midrule
c1 & 0.75 & M1RE\_INEQ\_02 & procedure & racines exactes et tableau complet \\
c2 & 0.75 & M1RE\_INEQ\_02 & automatisme & ensemble solution exact, bornes correctement ouvertes \\
c3 & 0.5 & M1RE\_INEQ\_02 & automatisme & ensemble complémentaire exact, bornes fermées incluses \\
\bottomrule\end{tabularx}
\minititle{Erreurs probables et codes}
\par\noindent\begin{tabularx}{\linewidth}{>{\ttfamily\small}p{26mm} X}
\toprule Code & Description \\ \midrule
CONCEPT & signe de $2-x$ traité comme celui de $x-2$ \\
METHODE & tableau de signes non construit alors qu'il est explicitement demandé, le signe étant annoncé sans justification \\
NOTATION & bornes incluses alors que l'inégalité est stricte \\
\bottomrule\end{tabularx}
\par\vspace{1mm}\small\textbf{Rôle pour la rentrée :} méthode reprise telle quelle pour le signe du trinôme
\par\vspace{2mm}
\itemhead{B2 --- \code{M1RE_FONC_02} --- noyau commun}{2}{3.75}
\begin{graybox}\small\textbf{Réponse attendue :} $5$ est l'image de $-2$ par $f$ ; $-2$ est un antécédent de $5$. Pour $0<x<1$, $x^2<x$.\end{graybox}
\minititle{Étapes significatives}
\begin{enumerate}\item $f(-2)=5$\item exemple : $x=0{,}5$ donne $x^2=0{,}25<0{,}5$\item argument : $x^2-x=x(x-1)$ avec $x>0$ et $x-1<0$, donc $x^2-x<0$\end{enumerate}
\minititle{Barème par critère --- la saisie se fait critère par critère}
\par\noindent\begin{tabularx}{\linewidth}{>{\ttfamily\scriptsize}p{9mm} >{\bfseries}p{9mm} >{\ttfamily\scriptsize}p{24mm} >{\scriptsize}p{18mm} X}
\toprule \scriptsize critère & Pts & \scriptsize compétence & \scriptsize preuve & Ce qui est exigé \\ \midrule
c1 & 1 & M1RE\_FONC\_02 & communication & vocabulaire correct dans les deux sens : image et antécédent \\
c2 & 0.6 & M1RE\_FREF\_02 & justification & argument valable pour tout $x$ de l'intervalle \\
c3 & 0.4 & M1RE\_FREF\_02 & procedure & exemple numérique cohérent \\
\bottomrule\end{tabularx}
\minititle{Erreurs probables et codes}
\par\noindent\begin{tabularx}{\linewidth}{>{\ttfamily\small}p{26mm} X}
\toprule Code & Description \\ \midrule
CONCEPT & image et antécédent échangés \\
CONCEPT & $x^2>x$ affirmé par généralisation du cas $x>1$ \\
JUSTIFICATION & un seul exemple donné, sans argument général \\
\bottomrule\end{tabularx}
\par\vspace{1mm}\small\textbf{Rôle pour la rentrée :} langage fonctionnel utilisé dans tous les chapitres de Première
\par\vspace{2mm}
\itemhead{B3 --- \code{M1RE_POURC_02} --- noyau commun}{2}{3.75}
\begin{graybox}\small\textbf{Réponse attendue :} $360$ TND puis $432$ TND ; coefficient global $0{,}9$, soit une baisse globale de $10\,\%$.\end{graybox}
\minititle{Étapes significatives}
\begin{enumerate}\item $480\times 0{,}75=360$\item $360\times 1{,}2=432$\item $0{,}75\times 1{,}2=0{,}9$ soit $-10\,\%$\end{enumerate}
\minititle{Barème par critère --- la saisie se fait critère par critère}
\par\noindent\begin{tabularx}{\linewidth}{>{\ttfamily\scriptsize}p{9mm} >{\bfseries}p{9mm} >{\ttfamily\scriptsize}p{24mm} >{\scriptsize}p{18mm} X}
\toprule \scriptsize critère & Pts & \scriptsize compétence & \scriptsize preuve & Ce qui est exigé \\ \midrule
c1 & 1 & M1RE\_POURC\_02 & procedure & les deux prix exacts \\
c2 & 1 & M1RE\_POURC\_01 & procedure & coefficient global et taux global exacts \\
\bottomrule\end{tabularx}
\minititle{Erreurs probables et codes}
\par\noindent\begin{tabularx}{\linewidth}{>{\ttfamily\small}p{26mm} X}
\toprule Code & Description \\ \midrule
METHODE & $-25\,\%$ et $+20\,\%$ additionnés en $-5\,\%$ \\
CONCEPT & coefficient d'une baisse de $25\,\%$ pris égal à $0{,}25$ \\
CALCUL & produit $0{,}75\times 1{,}2$ mal évalué alors que la méthode est correcte \\
\bottomrule\end{tabularx}
\par\vspace{1mm}\small\textbf{Rôle pour la rentrée :} base des suites géométriques et de l'évolution en pourcentage
\par\vspace{2mm}
\itemhead{B4 --- \code{M1RE_VECT_01} --- item individualisé}{2}{5.25}
\begin{graybox}\small\textbf{Réponse attendue :} $\vec{AB}(4\,;8)$ et $\vec{AC}(-2\,;-4)$ ; déterminant nul donc points alignés ; coefficient directeur $2$.\end{graybox}
\minititle{Étapes significatives}
\begin{enumerate}\item $\vec{AB}(7-3\,;6-(-2))=(4\,;8)$\item $\vec{AC}(1-3\,;-6-(-2))=(-2\,;-4)$\item $4\times(-4)-8\times(-2)=-16+16=0$\item $m=\frac{8}{4}=2$\end{enumerate}
\minititle{Barème par critère --- la saisie se fait critère par critère}
\par\noindent\begin{tabularx}{\linewidth}{>{\ttfamily\scriptsize}p{9mm} >{\bfseries}p{9mm} >{\ttfamily\scriptsize}p{24mm} >{\scriptsize}p{18mm} X}
\toprule \scriptsize critère & Pts & \scriptsize compétence & \scriptsize preuve & Ce qui est exigé \\ \midrule
c1 & 0.7 & M1RE\_VECT\_01 & automatisme & les deux vecteurs exacts \\
c2 & 0.7 & M1RE\_VECT\_02 & justification & alignement justifié par le calcul du déterminant \\
c3 & 0.6 & M1RE\_DROIT\_01 & procedure & coefficient directeur de $(AB)$ exact \\
\bottomrule\end{tabularx}
\minititle{Erreurs probables et codes}
\par\noindent\begin{tabularx}{\linewidth}{>{\ttfamily\small}p{26mm} X}
\toprule Code & Description \\ \midrule
CONCEPT & soustraction effectuée dans l'ordre $A - B$ \\
CALCUL & $-6-(-2)$ traité comme $-8$ \\
JUSTIFICATION & alignement affirmé sans déterminant \\
\bottomrule\end{tabularx}
\par\vspace{1mm}\small\textbf{Rôle pour la rentrée :} géométrie repérée, prérequis du produit scalaire
\par\vspace{2mm}
\itemhead{C1 --- \code{M1RE_DROIT_01} --- noyau commun}{4}{8.5}
\begin{graybox}\small\textbf{Réponse attendue :} $m=2$ et $y=2x$ ; $C$ appartient à $(AB)$ ; $P(N\cup C)=0{,}7$ et $P(\overline{N})=0{,}6$ ; $u_{n+1}=1{,}04\,u_n$ et $u_2\approx 195$ TND.\end{graybox}
\minititle{Étapes significatives}
\begin{enumerate}\item $m=\frac{10-2}{5-1}=2$ ; $y=2x+p$ avec $2=2\times 1+p$ donc $p=0$\item $2\times(-3)=-6$ : les coordonnées de $C$ vérifient l'équation (ou $\vec{AB}(4\,;8)$ et $\vec{AC}(-4\,;-8)$ sont colinéaires)\item $P(N)=0{,}4$, $P(C)=0{,}5$, $P(N\cap C)=0{,}2$ donc $P(N\cup C)=0{,}4+0{,}5-0{,}2=0{,}7$\item $P(\overline{N})=1-0{,}4=0{,}6$\item $u_{n+1}=1{,}04\,u_n$ ; $u_2=180\times 1{,}04^2=194{,}688\approx 195$\end{enumerate}
\minititle{Barème par critère --- la saisie se fait critère par critère}
\par\noindent\begin{tabularx}{\linewidth}{>{\ttfamily\scriptsize}p{9mm} >{\bfseries}p{9mm} >{\ttfamily\scriptsize}p{24mm} >{\scriptsize}p{18mm} X}
\toprule \scriptsize critère & Pts & \scriptsize compétence & \scriptsize preuve & Ce qui est exigé \\ \midrule
c1 & 1 & M1RE\_DROIT\_01 & transfert & coefficient directeur et équation réduite exacts \\
c2 & 1 & M1RE\_VECT\_02 & justification & appartenance justifiée par un calcul explicite \\
c3 & 1 & M1RE\_PROBA\_01 & transfert & les deux probabilités exactes, intersection correctement retirée \\
c4 & 1 & M1RE\_POURC\_02 & transfert & relation de récurrence correcte et $u_2$ exact \\
\bottomrule\end{tabularx}
\minititle{Erreurs probables et codes}
\par\noindent\begin{tabularx}{\linewidth}{>{\ttfamily\small}p{26mm} X}
\toprule Code & Description \\ \midrule
CALCUL & ordonnée à l'origine non déterminée, équation laissée en $y=2x+p$ \\
CONCEPT & $P(N\cup C)$ obtenue en additionnant sans retirer l'intersection \\
CONCEPT & $u_{n+1}=u_n+4$ : évolution additive au lieu de multiplicative \\
JUSTIFICATION & appartenance de $C$ affirmée sans calcul \\
\bottomrule\end{tabularx}
\par\vspace{1mm}\small\textbf{Rôle pour la rentrée :} quatre entrées de Première mobilisées sur une même situation ; la question 4 s'appuie sur la phase 4 de la séance
\par\vspace{2mm}
\itemhead{C2 --- \code{M1RE_NUM_01} --- item individualisé}{2}{4}
\begin{graybox}\small\textbf{Réponse attendue :} $3^{-1}=\frac13$. Non, l'affirmation est fausse : $1 < 2$ et pourtant $f(1)=1 > f(2)=0{,}5$. Un seul contre-exemple suffit à réfuter la croissance. (Deux valeurs ne démontrent en revanche pas la décroissance sur tout l'intervalle : cette propriété se démontre autrement.)\end{graybox}
\minititle{Étapes significatives}
\begin{enumerate}\item $3^{5}\times 3^{-2}=3^{3}$ puis $3^{3}\div 3^{4}=3^{-1}$\item deux valeurs bien choisies réfutent la croissance ; elles ne démontrent pas la décroissance globale\end{enumerate}
\minititle{Barème par critère --- la saisie se fait critère par critère}
\par\noindent\begin{tabularx}{\linewidth}{>{\ttfamily\scriptsize}p{9mm} >{\bfseries}p{9mm} >{\ttfamily\scriptsize}p{24mm} >{\scriptsize}p{18mm} X}
\toprule \scriptsize critère & Pts & \scriptsize compétence & \scriptsize preuve & Ce qui est exigé \\ \midrule
c1 & 1 & M1RE\_NUM\_01 & automatisme & puissance exacte \\
c2 & 1 & M1RE\_FREF\_01 & justification & affirmation réfutée par un contre-exemple explicite $f(1)>f(2)$ avec $1<2$ ; l'ajout « donc décroissante » n'est pas pénalisé, mais la décroissance globale n'est pas exigée \\
\bottomrule\end{tabularx}
\minititle{Erreurs probables et codes}
\par\noindent\begin{tabularx}{\linewidth}{>{\ttfamily\small}p{26mm} X}
\toprule Code & Description \\ \midrule
CONCEPT & exposants multipliés \\
CONCEPT & affirmation jugée exacte : la fonction inverse est déclarée croissante \\
JUSTIFICATION & réponse donnée sans contre-exemple chiffré \\
\bottomrule\end{tabularx}
\par\vspace{1mm}\small\textbf{Rôle pour la rentrée :} puissances et fonctions de référence, réactivées dès septembre
\par\vspace{2mm}
\newpage\sectiontitle{7. Correspondance diagnostic initial $\leftrightarrow$ évaluation finale}
\rowcolors{2}{SoftBlue}{white}
\par\noindent\begin{tabularx}{\linewidth}{>{\bfseries}p{9mm} >{\ttfamily\scriptsize}p{25mm} >{\ttfamily\scriptsize}p{22mm} >{\ttfamily\scriptsize}p{16mm} X}
\rowcolor{Navy}\color{white}Item & \color{white}\scriptsize skill\_id & \color{white}\scriptsize baseline & \color{white}\scriptsize comparaison & \color{white}Lecture \\
A1 & M1RE\_ALG\_02 & 1S\_TI\_Q02 & indicative\_skill\_comparison & confrontation indicative seulement : un énoncé parallèle existe au positionnement initial, mais la réponse de l'élève à cet énoncé n'a pas été conservée \\
A2 & M1RE\_INEQ\_01 & 1S\_TI\_Q03 & indicative\_skill\_comparison & confrontation indicative seulement : un énoncé parallèle existe au positionnement initial, mais la réponse de l'élève à cet énoncé n'a pas été conservée \\
A3 & M1RE\_FONC\_01 & 1S\_TI\_Q05 & indicative\_skill\_comparison & confrontation indicative seulement : un énoncé parallèle existe au positionnement initial, mais la réponse de l'élève à cet énoncé n'a pas été conservée \\
A4 & M1RE\_VECT\_01 & 1S\_TI\_Q09 & indicative\_skill\_comparison & confrontation indicative seulement : un énoncé parallèle existe au positionnement initial, mais la réponse de l'élève à cet énoncé n'a pas été conservée \\
A5 & M1RE\_FONC\_01 & 1S\_TI\_Q05 & indicative\_skill\_comparison & confrontation indicative seulement : un énoncé parallèle existe au positionnement initial, mais la réponse de l'élève à cet énoncé n'a pas été conservée \\
A6 & M1RE\_VECT\_02 & 1S\_TI\_Q10 & indicative\_skill\_comparison & confrontation indicative seulement : un énoncé parallèle existe au positionnement initial, mais la réponse de l'élève à cet énoncé n'a pas été conservée \\
B1 & M1RE\_INEQ\_02 & 1S\_TI\_Q04 & indicative\_skill\_comparison & confrontation indicative seulement : un énoncé parallèle existe au positionnement initial, mais la réponse de l'élève à cet énoncé n'a pas été conservée \\
B2 & M1RE\_FONC\_02 & 1S\_TI\_Q08 & indicative\_skill\_comparison & confrontation indicative seulement : un énoncé parallèle existe au positionnement initial, mais la réponse de l'élève à cet énoncé n'a pas été conservée \\
B3 & M1RE\_POURC\_02 & 1S\_TI\_Q14 & indicative\_skill\_comparison & confrontation indicative seulement : un énoncé parallèle existe au positionnement initial, mais la réponse de l'élève à cet énoncé n'a pas été conservée \\
B4 & M1RE\_VECT\_01 & 1S\_TI\_Q09 & indicative\_skill\_comparison & confrontation indicative seulement : un énoncé parallèle existe au positionnement initial, mais la réponse de l'élève à cet énoncé n'a pas été conservée \\
C1 & M1RE\_DROIT\_01 & 1S\_TI\_Q11+1S\_TI\_Q10... & indicative\_skill\_comparison & confrontation indicative seulement : un énoncé parallèle existe au positionnement initial, mais la réponse de l'élève à cet énoncé n'a pas été conservée \\
C2 & M1RE\_NUM\_01 & 1S\_TI\_Q17 & indicative\_skill\_comparison & confrontation indicative seulement : un énoncé parallèle existe au positionnement initial, mais la réponse de l'élève à cet énoncé n'a pas été conservée \\
\end{tabularx}
\begin{keybox}\textbf{Aucun écart chiffré ne sera calculé.} Les réponses de l'élève aux 18 questions du positionnement initial n'ont pas été conservées : le dossier n'en garde qu'un statut par domaine. Un statut qualitatif n'est pas une mesure : le convertir en score pour en soustraire un autre produirait un chiffre sans valeur. Le script d'analyse impose donc \code{mastery_delta = null} pour toutes les compétences, et se limite à une confrontation \emph{indicative} entre le statut initial et l'état final.\end{keybox}
\sectiontitle{8. Clés d'interprétation après correction}
\subsectiontitle{Échelle de maîtrise par compétence}
\par\noindent\begin{tabularx}{\linewidth}{>{\bfseries}p{10mm} X}
\toprule Niveau & Signification \\ \midrule
0 & non maîtrisé \\
1 & très fragile \\
2 & partiellement maîtrisé \\
3 & maîtrisé \\
4 & maîtrise solide / transfert réussi \\
\bottomrule\end{tabularx}
\subsectiontitle{Croiser réussite et certitude}
\par\noindent\begin{tabularx}{\linewidth}{>{\bfseries}p{40mm} X}
\toprule Situation & Conduite à tenir \\ \midrule
Juste, certitude 3--4 & acquis disponible : faire transférer sur une situation nouvelle. \\
Juste, certitude 1--2 & presque acquis : faire expliciter la procédure, puis réutiliser. \\
Faux, certitude 1--2 & notion à installer ou procédure à reprendre pas à pas. \\
Faux, certitude 3--4 & priorité absolue : confronter la conviction avant toute reprise technique. \\
\bottomrule\end{tabularx}
\subsectiontitle{Distinguer les niveaux d'erreur}
Une erreur de \textsc{CALCUL} isolée, non reproduite sur les items voisins de même compétence, ne vaut pas non-maîtrise conceptuelle. La chaîne à reconstituer est : \textbf{connaissance} $\rightarrow$ \textbf{choix de méthode} $\rightarrow$ \textbf{exécution} $\rightarrow$ \textbf{justification} $\rightarrow$ \textbf{contrôle}. Les codes d'erreur du barème situent l'écart sur cette chaîne.
\sectiontitle{9. Points d'observation propres à cet élève}
\begin{itemize}\item Indicateurs du dossier : aucune erreur de signe sur trois substitutions, quatre vecteurs exacts, colinéarité justifiée, pente calculée sans addition.\item Vérifier que les parenthèses sont écrites avant tout calcul lors d'une substitution négative.\item Le calcul littéral et les pourcentages étant solides, ne pas y consacrer de temps de consolidation.\end{itemize}
\subsectiontitle{Grille de saisie de la séance}
\par\noindent\begin{tabularx}{\linewidth}{>{\bfseries}p{26mm} X X X}
\toprule Phase & Aide maximale utilisée & Réussite autonome & Observation \\ \midrule
Phase 1 & & & \\[6mm]
Phase 2 & & & \\[6mm]
Phase 3 & & & \\[6mm]
Phase 4 & & & \\[6mm]
\bottomrule\end{tabularx}
\begin{warningbox}\textbf{Avant d'écrire quoi que ce soit sur les acquis de cet élève :} tant que l'évaluation n'est pas corrigée et analysée par \code{tools/analyze_s5.py}, aucune formulation du type « maîtrise désormais », « forte progression » ou « objectif atteint » ne doit être employée. Les formulations admises sont : « à vérifier lors de l'évaluation finale », « observé en séance, à confirmer », « hypothèse de consolidation ».\end{warningbox}
\end{document}
